\subsection*{Impactos económicos}
\addcontentsline{toc}{subsection}{Impactos económicos}

La implementación de la plataforma genera impactos económicos multidimensionales que trascienden el beneficio comercial inmediato, incidiendo tanto en el bienestar microeconómico de los usuarios como en el dinamismo macroeconómico del sector Fintech. A nivel microeconómico, el efecto más directo es la optimización del flujo de caja personal mediante el uso de presupuestos digitales y herramientas de simulación. Estas capacidades técnicas permiten reducir el endeudamiento improductivo y fomentar el ahorro sistemático, lo cual se traduce en un incremento del patrimonio neto de los usuarios a mediano y largo plazo. Para el segmento de trabajadores independientes e informales, que representa el 32,4\% de la población ocupada en Bogotá \parencite{DANE01}, el aplicativo actúa como un puente hacia la formalización, permitiendo la construcción de un historial financiero que facilita el acceso futuro a productos de crédito y servicios financieros formales.

Desde una perspectiva macroeconómica, el proyecto contribuye a la maduración y diversificación del ecosistema Fintech colombiano, el cual ha mostrado un crecimiento sostenido del 5,8\% anual \parencite{RevistaCIES01}. Al complementar la oferta de billeteras digitales y neobancos con una capa robusta de educación y planificación, se fortalece la resiliencia del sistema financiero. Según \textcite{BancadeOportunidades01}, la inclusión financiera efectiva no depende solo del acceso a cuentas, sino del uso productivo de las mismas; en este sentido, la plataforma incentiva la canalización del ahorro hacia vehículos de inversión formal, dinamizando el mercado de capitales local.

Finalmente, la integración de servicios basados en datos y arquitecturas abiertas potencia la eficiencia del mercado financiero. La adopción de modelos de \textit{Open Finance} permite una asignación de recursos más eficiente y reduce las asimetrías de información, factores que, según proyecciones de \textcite{Asobancaria01}, son determinantes para sostener el crecimiento del PIB sectorial en entornos de alta competitividad. De esta forma, el emprendimiento no solo busca la rentabilidad propia, sino que se posiciona como un catalizador de la estabilidad económica y el crecimiento del ecosistema digital en la región.